\begin{lemma}[Situation 7.1: shortening at the last exact edge]
Fix the Section 7 hierarchy counterexample and its supplied rung-local
\noderef{HierarchyMinimumCounterexampleContext}.  Assume
\(\noderef{HierarchyPathNumber}=1\) and suppose
\[
 \ell-1\in\noderef{HierarchyEligibleCommonMissingIndices}.
\]
Then these hypotheses contradict the supplied rung-local lexicographic
minimality and hence are impossible.
\end{lemma}
\begin{proof}
Expand the filter defining
\noderef{HierarchyEligibleCommonMissingIndices}.  It supplies a
\noderef{HierarchyReplacementPresentation} \(P\) and a color \(\alpha\)
missing under its coloring \(\mu_n\) at both
\[
 y_{\ell-1}=
   \noderef{HierarchyFullTailPrefix}(C,\ell-1).\mathit{lastVertex}
 \quad\hbox{and}\quad y_\ell=T.\mathit{lastVertex}.            \tag{4}
\]
The presentation supplies an actual corresponding ETT \(\mathcal E_P\)
on the same exact \(T\), maximum property, a base, a counterexample
\(C_P\), a full good hierarchy \(H_P\), and its transported minimum
context.

Apply \noderef{HierarchyCounterexampleTailAtLeastTwo} to that context.
The \noderef{SameGoodHierarchy} data identify the last starred stage with
the original one, so the resulting full-tail length is the same \(\ell\);
hence \(2\leq\ell\).  Put
\[
 B=\noderef{HierarchyFullTailPrefix}(C,\ell-1).
\]
Arithmetic from \(|A|+\ell=|T|\) gives \(|B|=|T|-1\).
The definition of \noderef{HierarchyFullTailPrefix} gives an exact
\noderef{TreeSequenceInitialSegment} \(B\preccurlyeq T\), retaining every
edge identity before the last one.

Let \(e_\ell=T.\mathit{edge}(\operatorname{Fin.last}|T|)\) be the actual
last edge.  Since
\(\noderef{HierarchyPathNumber}(H_P)=1\), its candidate condition at this
last exact index gives
\[
 G.\mathit{ends}(e_\ell)
       =\{y_{\ell-1},y_\ell\}.                                \tag{5}
\]
The equality of path numbers follows from the same stage and exact terminal
tree data in \noderef{SameGoodHierarchy}.  Because \(\ell\geq2\), this last
edge is not the exceptional first SE/RE seed.  Restrict the certified
terminal closure route to its final two prefixes.  Its last step is an
actual \noderef{TashkinovAugmentation}; hence
\[
 \mu_n(e_\ell)=\beta
\]
for some color \(\beta\), and \(\beta\) is missing at a vertex of the
predecessor of that step.  This predecessor is exactly \(B\), and the
witness occurs before the newly attached vertex \(y_{\ell-1}\); equivalently
\[
 \beta\in\overline\mu_n(T(y_{\ell-2})).                       \tag{6}
\]
The edge identity in this argument is \(e_\ell\), not merely an edge with
the endpoint pair in (5).

Define \(\sigma_n\) by changing only the value of this exact edge from
\(\beta\) to \(\alpha\).  Equation (4) and (5) say that \(\alpha\) is
missing at both endpoints, so no edge adjacent to \(e_\ell\) has color
\(\alpha\).  Thus \(\sigma_n\) is proper.  The augmentation gives
\(\mu_n(e_\ell)=\operatorname{some}\beta\), so this replacement neither
creates nor removes an uncolored edge; exactly \(e_0\) remains uncolored.
On every exact edge displayed in \(B\), \(\sigma_n=\mu_n\).
At vertices of \(B\) other than \(y_{\ell-1}\), all incident values are
unchanged.  At \(y_{\ell-1}\), \(\alpha\) ceases to be missing and
\(\beta\) becomes missing.

Apply \noderef{GoodHierarchyTerminalPrefixRepackage} to
\(\mathcal E_P\), its maximum property, base, hierarchy and goodness, and
the exact prefix \(B\).  The last starred stage is strictly inside \(B\)
because \(\ell\geq2\).  The theorem gives an ETT on \(B\) with the same
rung and complete certified series, maximum property, a rebuilt base and
hierarchy with the same divider count, every old nonterminal stage and
divider position, and a good hierarchy under \(\mu_n\).

We now verify goodness under \(\sigma_n\), rather than treating it as an
output of prefix repackaging.  Every hierarchy stage and starred stage is
contained in \(T(y_{\ell-2})\); hence its displayed edges and vertex
missing sets are unchanged.  The only changed boundary identity is
\(e_\ell\).  Its new color \(\alpha\) is not missing on any of these stages:
it is missing at \(y_{\ell-1}\) in the elementary penultimate prefix, so
elementarity excludes it at every earlier stage vertex.  Every Gamma color
is stage-missing by \noderef{GoodHierarchy}, and therefore \(\alpha\) is
not a Gamma color.  Replacing the old boundary color \(\beta\) by
\(\alpha\) consequently cannot create a forbidden Gamma-colored boundary
edge; if \(\beta\) was controlled, removing it only helps closedness.  This
preserves the positive-stage minus-closed and PE parallel-base boundary
clauses.  For the terminal
avoidance condition, only the first-missing positions of \(\alpha\) and
\(\beta\) can change.  The color \(\beta\) already has the earlier witness
in (6), so its least position is unchanged.  If \(\alpha\) has an earlier
missing witness, that witness is unchanged; if its only missing witness in
\(B\) was \(y_{\ell-1}\), after recoloring it has no witness, and the
totalized fallback in \noderef{HierarchyFirstMissingPosition} is the same
last index \(y_{\ell-1}\).  Thus every stored color has the same first
missing position.  The segment-color sets in
\noderef{HierarchySegmentColors} consequently use the same exact bounds
and the same displayed edge values.

The Gamma pairs, their cardinalities and disjointness, all stage-missing
sets, every positive-stage minus-closed clause, and the PE parallel-base
clause therefore copy from the repackaged good hierarchy.  This checks all
fields of \noderef{GoodHierarchy} under \(\sigma_n\).  It is important here
that no equality of the full terminal missing-color set was asserted:
the controlled exchange at \(y_{\ell-1}\) is handled by (6) and by the
first-position fallback.

The new vertex set \(V(B)\) is not elementary under \(\sigma_n\):
equation (6) gives an earlier vertex missing \(\beta\), and the last-edge
recoloring makes \(\beta\) missing at \(y_{\ell-1}\); the two vertices are
distinct by the injective displayed vertex map in
\noderef{TreeSequence}.  Together with positivity, maximum property, the
rebuilt hierarchy, and the goodness just proved, this constructs a
\noderef{HierarchyCounterexample} \(C'\) on the exact shorter tree \(B\).

The ETT of \(C'\) retains the complete series of \(\mathcal E_P\), so the
rung numbers agree, and prefix repackaging retained the divider count.
Its path number is one: the old candidate \(1\) still satisfies the
consecutive-edge condition on every exact tail edge remaining in \(B\),
and \noderef{HierarchyCounterexamplePathNumberBounds} excludes zero.
The last ordinary hierarchy stage is unchanged, while
\(|B|=|T|-1\).  Therefore
\[
 \noderef{HierarchyCounterexampleMeasure}(C')
   =\bigl(1,\ |T|-1-|T_{n,q}|\bigr)
   <_{\rm lex}
    \bigl(1,\ |T|-|T_{n,q}|\bigr)
   =\noderef{HierarchyCounterexampleMeasure}(C_P)
\]
in the exact order \noderef{HierarchyCounterexampleLess}.  This contradicts
the lexicographic-minimality field of the transported minimum context in
\(P\).  Hence the assumed eligible penultimate index is impossible.
\end{proof}
